\documentclass[11pt,a4paper]{article}
\usepackage[margin=22mm]{geometry}
\usepackage{amsmath,amssymb,amsthm,booktabs,microtype}
\usepackage[T1]{fontenc}\usepackage{lmodern}
\usepackage[colorlinks=true,linkcolor=blue,urlcolor=blue]{hyperref}
\newcommand{\FS}{\operatorname{FS}}\newcommand{\Z}{\mathbb Z}
\newcommand{\T}{\mathcal T}
\newtheorem{theorem}{Theorem}[section]\newtheorem{lemma}[theorem]{Lemma}\newtheorem{proposition}[theorem]{Proposition}
\title{Subset Sums in a Fourfold Sumset\\\large Progress toward cardinality 89:\\tighter budgets and an exact relaxation obstruction}
\author{Sum-of-Subset research project}\date{18 October 2026}
\begin{document}\maketitle
\begin{abstract}
The independently audited global bound remains $N\le89$.
For a hypothetical 89-element configuration, we improve the low-hole R-band
budget from $329-I$ to $313-I$, removing its entire four-empty family.
Intrinsic normalization repairs sixteen failed sufficient inverse directions;
early prefix retention avoids a circular use of lattice peeling.
Further refinements count core overlap, absorbed peeled points and external
pair labels in interval unions. The supplied partial certificates exclude
3,651 structural groups, representing 95,443 integer root systems.
However, an exact rational point satisfies all 1,619 root inequalities and
one branch of every one of the 682 available dichotomies in an unresolved
balanced mode. Thus the present relaxation cannot prove $N\le88$.
The result is verified research progress, not a new global upper bound.
\end{abstract}

\section{Scope and endpoint constants}
For a finite real set $A$, let $\FS(A)$ include the empty subset sum, and let
$4A$ permit repeated summands. The quantity $N$ is the largest size with
$\FS(A)\subseteq4A$. We continue the independently audited chain ending at
$N\le89$ \cite{previous}; the lower construction has 71 elements. Any feasible
89-element set contains zero, because adjoining it otherwise gives
90 elements. Feasibility is not assumed to pass to arbitrary subsets.

The corrected integralization and peeling reduction \cite{foundation,previous}
give $D\supset D_1\supset E=P\cup\{0\}\cup(-Q)$, with
\[
 d=81,\quad e=73,\quad a+b=m=72,\quad I=(a-36)^2,
 \quad\Lambda=1+T_a+T_b=1333+I,
\]
where $T_r=r(r+1)/2$. Set
\[
 \Omega=22d+97-\Lambda=546-I,\qquad
 C=2m+4e+14-\Lambda=-883-I.
\]
For a positive $r$-set $X$, write
$\mu_X=|\FS(X)|-1-T_r$ and $\rho_X=\mu_X+2-2\delta_0(X)$,
where $\delta_0$ counts defective diagonals in the distinct-pair sumset of
$X\cup\{0\}$. The maximum is defective, hence $\rho_X\le\mu_X$.
The inherited bounds are $\rho_X\ge-8$ universally, $\rho_X\ge0$ for
$r\ge14$, and, for non-progression support,
\[
 \mu_X\ge r-3,\qquad \rho_X\ge f(r),\qquad
 f(r)=\begin{cases}r-14&14\le r<17,\\r-5&r\ge17.\end{cases}
\]
The exact budget is
\begin{equation}\label{budget}
 \rho_P+\rho_Q+2\vartheta+c+\T=\Omega,
 \quad c=2|2D|+18d+85-|\FS(D)|,
 \quad\T=\kappa+\ell_1+\ell_2+g,
\end{equation}
with $\kappa=|\FS(P)-\FS(Q)|-|\FS(P)|-|\FS(Q)|+1$.
The slack quantities $\vartheta,c,\kappa,\ell_1,\ell_2,g$ are nonnegative. Only after proving the common lattice for
$D_1,D$ may we use
\begin{equation}\label{sharp}
 w\le4z,\quad \T\ge\kappa+152-3z,\quad
 \T\ge\max(51,z,152-2z),\qquad z=h_P+h_Q-\kappa.
\end{equation}

Here $w$ denotes the number of hull holes of $\FS(D)$. In particular, $z=0$ makes its two target sides complete integer intervals.

\section{Primitive normalization repairs the inverse test}
The deficient-sign check is $1+T_{73}=2702>24\cdot89+77=2213$.
The paired bootstrap has $q^2\le2752$, even $q\le52$, and hence
$a,b\ge10$. Sizes below 13 have $\Omega<-8$ and are impossible.
For sizes $13,59$, $\Omega=17$. A non-progression larger side costs at
least $54-8>17$. Otherwise its progression support forces one residue and
no long gap, since $T_{59}>25$; then $z=0$. The inner comparison, with the
mandatory small-prefix correction, is bounded by $C+3(\Omega+8)+16<0$;
the outer comparison is also strict. Formula~\eqref{sharp} would give
$\T\ge152>25$, a contradiction.

For $a,b\ge14$, the larger-side strict inverse margin is at least 118.
When both supports are non-progressions, the sufficient smaller-side test is
\begin{equation}\label{strict}
 \mu_P+\rho_P+\rho_Q>\Omega-T_a+2,
 \qquad a\le b.
\end{equation}
Its floor margin can be as small as $-21$, and fails exactly for
$20\le a\le27$. We treat actual failure of \eqref{strict}, not just failure
of its conservative floor test.

\begin{lemma}[Primitive form of the scheduled hole bound]\label{primitive}
The variable-threshold persistent-hole theorem in \cite{schedule} remains
valid for a primitive positive integer set without the upper hypothesis
$h_X<|\FS(X)|-1$.
\end{lemma}
\begin{proof}
That upper hypothesis is used only to deduce primitivity. With primitivity
given, $h_X>0$ excludes progression support. For a proposed lower defect
$B$, the same schedule checks the following inequalities:
\begin{gather*}
 18\le L\le r,\quad 2\beta_i<z,\quad 2\beta_i+4\le i,\quad\beta_i\ge2,\\
 B\le T_{L-1}-r-1,\quad B\le T_i+i-6\beta_i-r-7,
 \quad B\le L-6+\sum_{i=L}^r(\beta_i-1).
\end{gather*}
Assuming $\rho_X<B$, the first escape excludes a late progression prefix;
the recurrence $\rho_i\ge\rho_{i-1}+y_i-2$ then either gives the last
inequality's contradiction or supplies a small-increment prefix.
The small-increment inverse writes that prefix as $tY$, with
$\max Y\le i+\beta_i$. A later first coefficient off $t\Z$ produces
disjoint translates and violates the same escape bound. Primitivity gives
$t=1$. The central interval theorem supplies only endpoint holes, at most
$2\beta_i$ in total. The second escape excludes every later large jump;
all subsequent central intervals join and preserve those endpoint holes.
This contradicts $h_X\ge z>2\beta_i$. None of these steps uses an upper
bound on $h_X$.
\end{proof}

If \eqref{strict} fails, put
\[
 R=\Omega-T_a+2-(a-3)-(a-5),\quad K=\Omega-62,
 \quad U_{\min}=T_a+a-2,\quad U_{\max}=K+3,\quad V_{\min}=T_b+b-2.
\]
Then $\rho_Q\le R$, $\kappa\le K$, and
$U_{\min}\le|\FS(P)|\le U_{\max}$, $|\FS(Q)|\ge V_{\min}$.
Normalize the initially integralized $Q$ by its own gcd. Apply
Lemma~\ref{primitive} with $z=7$, $L=18$, and constant threshold $\beta_i=3$:
these schedules give $B=2b-22$, each a defect lower
bound greater than $R$. Thus $h_Q\le H=6$ on this intrinsic lattice.
The values of $R$ for $a=20,\ldots,27$ are
$50,58,63,65,64,60,53,43$.

Here is the complete alignment argument. Let $U=\FS(P)$, $V=\FS(Q)$,
$u=|U|$, $v=|V|$, and $A,B$ their diameters. If $U$ occupied $s$ residue
classes modulo this lattice, residue-wise sumset counting would give
$\kappa\ge(s-1)(v-1)$. Since $V_{\min}-1-K>0$, there is only one
class; zero belongs to it. Thus $P$ is also integral. If $A>B$, the two
endpoint translates $V$ and $A+V$ are disjoint, forcing
$v\le u-1+\kappa$, contrary to $V_{\min}-U_{\max}-K+1>0$.
Hence $A\le B$. The two endpoint hulls cover the sum hull, so
\[
 |U+V|\ge A+B+1-2h_Q=A+v-h_Q,
 \qquad h_P\le\kappa+H.
\]
The further margins $V_{\min}-K-H-2>0$ and $U_{\min}-H-2>0$
give $\max(A,B)\le u+v-3$. Stanchescu's Theorem 4
\cite{stanchescu} therefore gives the exact bounds
$h_P,h_Q\le\kappa$. Its diameter hypothesis includes the equal-diameter
case and requires no extra individual gcd assumption.
The sum hull also gives $z\ge0$, while $z\le\min(h_P,h_Q)\le6$.

The four minimum margins, in the order used above, are
\[674,\qquad270,\qquad667,\qquad220.\]
The pre-lattice inner bound is now
$C+\Omega+2H<0$, and the outer margin is at least 473. Thus lattice
propagation is complete before using \eqref{sharp}. It gives
\[
 \kappa\le\Omega-62-152+18<U_{\min}-1,
\]
with strict surplus at least 115. Consequently both side-hole bounds are
strict, and both sides are primitive. This closes all 16 oriented repairs.
The derivation uses the general short-diameter method recorded in the
archived structural notes \cite{archive}; no archived endpoint certificate
or unreviewed endpoint conclusion is imported.

\section{Early retention closes the remaining lattice comparison}
Write $F_0=|\FS(E)|$ and $F_1=|\FS(D_1)|$.
For the actual strict-inverse branches, the common core lattice and strict
side-hole bounds already hold. The inherited comparisons are
\[
 F_1-2F_0\le\min\{C+\Omega+\rho-\mu+2z,
                         C+3\Omega-\rho-\mu\},
 \quad\rho=\rho_P+\rho_Q,\quad\mu=\mu_P+\mu_Q.
\]
The first is at most $-337-2I+2z$, hence strict for $z\le168+I$.
For $z\ge Z=169+I$, the basic defect and support floors settle most
orientations. On every remaining side $X$ of size $r$, the pre-lattice
budget itself gives
\[
 \rho_X\le B=\Omega-f(72-r)-Z<T_{r-1}-r-1.
\]
The exact verifier checks this strict escape for both sides in every
remaining orientation. It uses neither $G=152$ nor a sharp peeling bound.

The final prefix translates must intersect. Indeed disjointness would give
$y_r\ge1+T_{r-1}-r$ and, from $\rho_{r-1}\ge0$,
$\rho_r\ge T_{r-1}-r-1>B$. If $Y=X\setminus\{\max X\}$,
intersecting hulls imply $h_X\le2h_Y$. Maximum removal therefore gives
\[
 \max X\le r+\rho_X+h_X+2-\rho_Y-h_Y
          \le r+\rho_X+h_X/2+2.
\]
Insert this on both sides in the original inner comparison:
\[
 F_1-2F_0\le C+2\rho-\mu+z+\ell_2
 \le C+\Omega+\rho-\mu\le-337-2I<0.
\]
The middle step uses $\rho+\kappa+\ell_2\le\Omega$ and $z\le\kappa$.
This proves the missing integrality without a circular lattice assumption.
For a side of size 14, use the universal prefix floor $-8$ and add 16
to the conservative comparison; $C+3\Omega+16<0$ still holds.
The outer margin is always $473+2I>0$.

\section{Progression supports and small surviving sides}
The progression-support repair of \cite{schedule,previous} is repeated
with the new constants. An AP side of size at least 36 has support length
at least $T_{36}=666>\Omega$. Residue counting and the exact formula
\[
 |W+[0,u-1]_{\Z}|=|W|+u-1+
       \sum_i\min(w_{i+1}-w_i-1,u-1)
\]
force one lattice residue and no clipped gap, giving $z=0$.
For a smaller AP side, $\Omega<2T_a$ and $\Omega-T_a\le102$.
A unique clipped gap or a sole off-lattice coefficient makes the last
support translates on the other side disjoint. Here the clipping or
residue cost already gives $\kappa\ge T_a$; retaining it yields the
strict contradiction
\[
 T_a+T_{b-2}+2-2b>\Omega.
\]
This extra $T_a$ must not be dropped at the balanced end of this argument.

The remaining two-residue case has at least two off-lattice coefficients
and residue order two. In half-step normalization $U=2[0,u-1]$, let
$w$ be the least odd coefficient of $Q$. The inherited fiber identities are
\[
 \kappa=u-1+h_0+h_1,\quad h_P=u-1,\quad
 h_Q=\kappa-u+w,\quad z=w-1,
 \quad\mu_Q+\kappa=\sum Q+u-w-T_b.
\]
If there are $s\ge2$ odd coefficients, distinctness gives
$\sum Q\ge(b-s)(b-s+1)+sw+s(s-1)$.
The exact mass comparison first excludes $w>u$. If $w\ge15$, its quadratic
minimum gives $\sum Q-w-T_b\ge7b-289/8$, whereas the budget gives
at most $2b+99$. This is a strict contradiction throughout the relevant
range. Thus $w\le13$ and $z\le12$.
The pre-lattice inner comparison is
\[
 F_1-2F_0\le2483+2a-4T_a-2T_b\le-1021.
\]
The outer comparison is strict, so \eqref{sharp} now applies:
$\T\ge T_a+116$ contradicts $\Omega-T_a\le102$, with surplus 14.
Surviving AP cases have $z=0$; their side-hole bounds are strict after
$\kappa\le\Omega-152$. Finally, for positive holes and a side below 17,
apply the scheduled bound only to the larger side. The scalar budget
excludes every such candidate. All surviving positive-hole sides have
size at least 18.

\section{Post-lattice retention and exclusion of a single hole}
The prefix-retention lemma of \cite{previous} extends to $r-q\ge14$:
under $\rho_X\le B<T_{r-q}-r-1$, each last $q$ prefix step intersects,
primitivity and non-progression persist, and
\[
 h_{X_{r-j}}\ge2^{-j}h_X,\quad
 \rho(X_{r-j})\ge f(r-j),\quad\mu(X_{r-j})\ge r-j-3.
\]
The same backwards proof works because the escape uses only the
nonnegative prefix defect from size 14. At sizes 14--16 the sharp
$r-j-5$ defect floor is replaced by the actual $f(r-j)$.
The support-excess floor remains valid there. For the endpoint,
$q=2$ is licensed at every surviving positive-hole state, with minimum
surpluses 39 and 15 for the first and second steps. The smallest first
prefix has size 17, so the span still uses the sharp first-prefix floor.

\begin{lemma}\label{onehole}
If an integer subset-sum support has exactly one hull hole, its final
prefix translates are disjoint.
\end{lemma}
\begin{proof}
Take the first coefficient $x$ that fails interval coverage, and let $L$
be the total of its predecessors. Their support is $[0,L]$, and
$x\ge L+2$. The missing value $L+1$ cannot be filled by later positive
coefficients. Since complement symmetry fixes the unique hole, the final
total is $2L+2$. Any later coefficient would exceed $x$ and make this total
larger than $2L+2$. Thus $x$ is final, equals $L+2$, and its translate
is disjoint from $[0,L]$.
\end{proof}
For $Y=X_{r-1}$, the second retained step intersects. Since $Y$ is
primitive and non-progression, it has positive holes; Lemma~\ref{onehole}
strengthens this to $h_Y\ge2$.

Let $\delta=c-81j$ for $j$ empty ordinary units. The ordinary retained
span and this additional two-hole floor give, respectively,
\begin{align}
 R_p+R_q+\delta-3z&\le410-I-81j,\label{spanhalf}\\
 R_p+R_q+\delta-4z&\le406-I-81j.\label{span2}
\end{align}
Indeed maximum removal and the budget give
$R_p+R_q+c\le470-I+4z-(\rho_{P'}+\rho_{Q'})-(h_{P'}+h_{Q'})$;
the first subtraction is at least $m-12=60$, and the second is at least
both $z$ and 4. The second row is new and is stronger when $z<4$.

For one side, with coordinates $(R,x_1,x_2,x_3,M)$ and total mass $S$,
the former half-hole mass rows are retained. Two additional necessary rows are
\begin{align}
 M+4R-S&\le-T_{r-1}-6-(r-4)-2,\\
 R+x_1+x_2+x_3+M-S&\le-T_{r-1}-(r-4)-2.
\end{align}
Both follow from $\sum X_{r-1}=T_{r-1}+\mu_{r-1}+h_{r-1}$.
The two-prefix mass row still uses $\mu_{r-2}\ge r-5$; at $r=18$
this is a support-excess statement for size 16, not the unproved sharp
defect statement for size 16.

\section{Core overlap across the original covering groups}
Two ordinary centers from different original groups can coincide or be very
close.
For example, the centers labelled $p5$ and $p8$ are equal when
$p_1-q_2=2p_2$. Grouping only consecutive own or cross centers misses this
possible repeated coverage.

Here is a uniform bound which does not assume coincidence. Let two contributing
ordinary centers $u,v$ belong to the same target side. Their full translates
have size $d$, and their core translates have size $e$. Put
$\lambda_u=|(u+D)\setminus F|$, $\lambda_v=|(v+D)\setminus F|$ for that
target $F$. The two translated cores fit in an integer interval of at most
$R_p+R_q+|u-v|+1$ points. Adding the $2(d-e)$ labelled tail points gives a
container for both full translates. This container includes all the points
counted by either member loss. Thus their common deficit satisfies
\begin{equation}\label{corepair}
 2d-|((u+D)\cup(v+D))\cap F|
 \ge 2e-R_p-R_q-|u-v|-1+\max(\lambda_u,\lambda_v).
\end{equation}
Here and below the union is intersected with $F$ as a whole.
The centers are integral because each contributes to the integer target and
$D$ is integral; no unproved integrality of all original outer elements is needed.

Let $L_{\rm rest}$ be the sum of the other ordinary member losses and let
$L_{\rm fixed}$ be the fixed-label loss already separately justified.
By bounding the full covering union with this pair union and the other
labelled sets, its remaining deficit obeys
\[
 \delta\ge 2e-R_p-R_q-|u-v|-1+
          \max(\lambda_u,\lambda_v)+L_{\rm rest}+L_{\rm fixed}.
\]
Each model loss variable is at most its corresponding actual loss, so replacing
actual losses by those variables preserves necessity. For every unordered
pair on the same target side, write $u-v$ in the eight positive outer-gap
coordinates. If all coefficients have one sign, that order is proved and
the two member-loss rows are inserted directly in the root model. Otherwise
the checker supplies both center orders $u-v\ge0$ and $u-v\le0$. In each branch it inserts the two rows choosing
one member loss at a time. Their union is exhaustive even at equality.
These pair bounds are alternatives; they are never summed as independent
savings for overlapping pairs.
In the matrix convention $Ax\le b$, the hull-width and distance coefficients
are negative:
\[
 -R_p-R_q-|u-v|-\delta+L_{\rm rest}+\lambda_u
 \le 1-2e-L_{\rm fixed}.
\]
This sign is also checked by evaluating the generated rows on actual sets.
 The binary certificate tree selects a valid
bound while preserving both order branches.

\paragraph{Absorbing peeled points into the same container.}
The pair estimate admits a further strengthening without adding score
indicators. Orient the centers as $u>v$ and put $\Delta=u-v$.
On the target side $s$, let $X_1<X_2<X_3<X_4$ be the four positive
peeled coordinates outside the core, whose positive endpoint is $R_s$.
If $X_k\le\Delta+R_s$, the first $k$ positive peeled points of $v+D$
already lie inside the two-core hull. They should not be counted again
among its $2(d-e)$ additional labelled points. The right side of
\eqref{corepair} therefore increases by $k$. Similarly, if $Y_k\le
\Delta+R_t$, the first $k$ negative peeled points of $u+D$ are already
inside the hull, where $t$ is the opposite side.

For each ordered pair and each such coordinate $X_k$, the verifier uses
exactly the two alternatives
\[
 \Delta+R_s-X_k\ge0,\qquad \Delta+R_s-X_k\le-1.
\]
These are exhaustive because contributing centers and all peeled-layer
coordinates are integers. In the first branch $X_k>R_s$ already implies
$\Delta>0$, so both the orientation and the strengthened pair inequality
are justified. The second branch makes no extra gain claim. The negative
peeled coordinates use the same argument with $R_t$. Gains from different
thresholds are not added; each accepted row is one complete container bound.
The general proof is supplemented by 199,624 checks of these generated
matrix rows on actual finite sets.

\paragraph{Joint absorption and disjoint pairs.}
The positive peeled points of the lower translate and the negative peeled
points of the higher translate are distinct families of labels. Thus their
absorption counts $A_+,A_-\in\{0,\ldots,4\}$ may be added within the
same container estimate. Alternative thresholds for the same family provide
a maximum count, not a sum. For the two distinguished pairs $(p5,p8)$ and
$(q5,q8)$ in the exceptional A-mode, both center orders carry these two
bounded counts. The threshold $\Delta+R_s-X_k\ge0$ forces $A_+\ge k$;
its complement forces $A_+\le k-1$, and the negative family is analogous.

Let $I_s$ be the actual intersection size inside the target interval for
one of these pairs. Since its pair deficit equals $\lambda_u+\lambda_v+I_s$,
\eqref{corepair} with both absorption counts gives
\[
 I_s\ge2e-R_p-R_q-|u-v|-1+A_++A_- -\min(\lambda_u,\lambda_v).
\]
Each member yields one necessary linear lower bound for $I_s$. In this
zero-hole mode the lower and upper tail counts sum to the exact member loss.
The two selected pairs use disjoint ordinary labels, so their intersections
can both be charged:
\[
 \delta\ge L_{\rm fixed}+\sum_{\rm ordinary}\lambda_v+I_p+I_q.
\]
This follows by bounding the full union by the two pair unions and all
remaining individual sets. It does not add deficits from overlapping pairs.
Actual-set checks include 27,108 joint-absorption rows and 2,259 disjoint-pair
allocations.

\paragraph{Two-way extreme-point quantiles in the zero-hole band.}
Write $Y^{(i)}$ for the $i$th largest negative magnitude and $X^{(i)}$
for the $i$th largest positive coordinate of $D$. The nested coordinates
provide these values for $1\le i\le6$. For a contributing center $c$ and
target interval $[1,S]$, use the actual lower and upper tail counts $l,u$.
The following are exact exhaustive alternatives:
\begin{align*}
 &(c-Y^{(i)}\le0,\ l\ge i)
       &&\text{or }(c-Y^{(i)}\ge1,\ l\le i-1),\\
 &(c+X^{(i)}-S\ge1,\ u\ge i)
       &&\text{or }(c+X^{(i)}-S\le0,\ u\le i-1).
\end{align*}
For example, if $c-Y^{(i)}\ge1$, all but at most the $i-1$ smaller
translated points are positive. This proves the first upper bound; the
other three statements follow by the same ordering argument. All thresholds
are integral. The old one-way alternatives had omitted the companion
position condition and the opposite upper count bound. Adding both retains
the same number of binary alternatives but excludes fictitious tail losses.
These strengthened alternatives are used only when $z=0$, so no unproved
exact-loss transport across positive-hole bands is required. Positive-hole
models retain their earlier one-way disjunctions. There are 108,432 finite
checks of the two-way statement on actual translated sets.

\paragraph{Count only outer-pair labels outside the central two-sum hull.}
Let the four positive peeled coordinates be $X_1<X_2<X_3<X_4=M_p$,
and similarly $Y_1<Y_2<Y_3<Y_4=M_q$. Every sum in $2D$ outside
\[
 H=[-M_q-R_q,\ M_p+R_p]_{\Z}
\]
uses two positive peeled coordinates, or two negative peeled coordinates.
Mixed-sign sums and any same-sign pair containing a core point lie in $H$.
Let $J_p$ count the ten unordered positive pair labels $i\le j$ for which
$X_i+X_j>M_p+R_p$, and define $J_q$ analogously. Labels may coincide;
counting them separately is a valid upper bound. Thus
\[
 |2D|\le M_p+M_q+R_p+R_q+1+J_p+J_q,\qquad 4\le J_p,J_q\le10.
\]
The four labels containing $X_4$ or $Y_4$ always lie outside the respective
endpoint. Only the six labels among the first three coordinates need tests.

For $1\le i\le j\le3$, let $A_{ij}$ count pairs $(k,l)$ with
$1\le k\le l\le3$, $k\le i$, $l\le j$, and let $B_{ij}$ count those
with $k\ge i$, $l\ge j$. If $X_i+X_j\le M_p+R_p$, at least $A_{ij}$
labels are inside, so $J_p\le10-A_{ij}$. Otherwise integrality gives
$X_i+X_j\ge M_p+R_p+1$ and $J_p\ge4+B_{ij}$. These two branches are
exhaustive; the negative side is identical. They use two bounded count
variables and twelve binary thresholds in total. In particular, if
$2X_3\le M_p+R_p$ and $2Y_3\le M_q+R_q$, the additive constant 21 in
the old hull estimate becomes 9. A further 27,108 finite threshold checks
and the corresponding actual two-sum sets supplement this general proof.

\paragraph{Deduplicate external pair labels in an interval union.}
The two exceptional four-empty modes also use six cuts
$(i,h)\in\{(3,1),(1,3),(3,2),(2,3),(2,2),(3,3)\}$.
Put $X_0=R_p$, $Y_0=R_q$ and use the four peeled coordinates as above.
After retaining the hull $C=[-Y_h,X_i]_{\Z}$, the external points are
$T=\{X_{i+1},\ldots,X_4,-Y_{h+1},\ldots,-Y_4\}$, of size $k=8-i-h$.
All non-external-pair sums lie in
\[
 U=2C\ \cup\ \bigcup_{t\in T}(t+C).
\]
With $W=X_i+Y_h$, the integer interval union has size
\[
 |U|=2W+1+\sum_{j=i+1}^4\min(W+1,X_j-X_{j-1})
             +\sum_{j=h+1}^4\min(W+1,Y_j-Y_{j-1}).
\]
Indeed the positive translated intervals extend the right end in increasing
order, each adding the smaller of its shift and its length; the negative
intervals do the same on the left. Their added parts are on opposite sides
of the central interval. Consequently
\[
 |2D|\le |U|+\frac{k(k+1)}2-
     \#\{\text{unordered external-pair labels whose sum belongs to }U\}.
\]
Every deducted label was counted once by the pair-label bound and already
belongs to the interval union; coinciding labels cause no difficulty.

Membership in each interval is two linear inequalities. The verifier
reconstructs all 45 pair labels and 78 nontrivial integer membership guards;
known signs are proved from positive nested gaps. A guard $f\ge0$ is
encoded by the exhaustive alternatives $f\ge0$ and $f\le-1$.
When both endpoint guards of any interval hold, its pair-membership variable
is forced to one. The minimum terms in $|U|$ use auxiliary variables bounded
above by both entries, with the actual minima providing a necessary model
assignment. There are 10,800 actual-set checks of the complete envelope.
For $(i,h)=(3,3)$, this removes the extra count for $M_p-M_q$ whenever
that mixed sum already lies in $[-2Y_3,2X_3]$.

\section{Covering domain and verification boundary}
The model reconstructs all ordinary contributor modes and integer row-lattice
witnesses from \cite{previous}. Quantile thresholds now are global $81/80$,
core $77/76$, and leading $161$, $153$, $151$. A group's actual cap decides
which rows are licensed; the high-hole cap 184 does not uniformly license
the basic leading rows. The cardinality row has right side $1542-81j$.
Low holes are partitioned as follows:
\begin{center}\begin{tabular}{crrr}\toprule
Band&$z$&combined defect floor&cap before $I+81j$\\\midrule
Z&0&0&394\\
P&1--2&62&336\\
Q&3--4&74&328\\
R&5--14&93&313\\\bottomrule
\end{tabular}\end{center}
The P and Q caps exceed $4d=324$, so both four-empty families must
be retained. For R, keeping the same $z$ in the defect estimate proves
\[
 \rho_P+\rho_Q-2z\ge81\qquad(17\le a\le55,\ 5\le z\le14).
\]
The 390 integer cases each carry an effective-hole witness, an integer
inverse witness and a scheduled-hole witness; use the largest proved side
floor and then subtract this actual $2z$. Combining with
$\T\ge152-2z$ yields $c\le394-I-81=313-I<4d$.
Taking the defect minimum and hole maximum independently had only given
$329-I$. The coupled comparison removes the entire R4 family before covering
models are built. High-hole states allow two empty units; these are also
present. Each low band transports $z$ to its upper endpoint while increasing
both target masses by the same amount. Every model row is checked for
preservation under that transport. High groups keep a common hole variable,
with all three coupled budget caps, and every actual state's rows imply
the target group's rows.

The full-translate losses and inherited ordinary disjunctions remain as in
\cite{previous}. The core-overlap rows, absorbed-tail dichotomies, two-way quantiles and
two outer-pair count variables strengthen this endpoint model. The inherited
110-score indicator extension, order rows and sign clauses are retained
only in the two exceptional zero-hole four-empty modes; all of their
hypotheses and complete score dichotomies remain as in \cite{previous}.
The eight already justified high fixed-label losses
remain charged where their hypotheses hold. Every tree has both children;
only exact nonnegative integer multipliers with $wA\ge0$ and $wb<0$ certify
a leaf. The standard-library checker reconstructs the complete root universe.
Numerical optimization discovers weights but has no role in acceptance.

\section{Exact obstruction and completed verification}
The unresolved root is
\path{Z4:A:pair1,pair2,p0,q0:-:-}, with $a=b=36$.
The file \path{data/upper88_relaxation_witness.json} contains a rational
assignment to all 450 model variables and a branch choice for every available
dichotomy. The standard-library checker verifies nonnegativity, all 1,619
root rows and all 1,420 rows in the 682 selected alternatives using
\texttt{fractions.Fraction}. There is no floating-point tolerance.

\begin{proposition}
The current root relaxation and its listed dichotomies cannot exclude this
root by any complete tree of their nonnegative linear certificates.
\end{proposition}
\begin{proof}
The same rational point satisfies the root and at least one alternative of
every listed dichotomy. Follow a satisfied child at each node of any complete
finite tree. The point satisfies every inequality at the resulting leaf,
so nonnegative linear combination cannot give a strict contradiction there.
\end{proof}
This point is a feasible point of a necessary relaxation. It is not asserted
to represent an actual feasible set for the original subset-sum problem.
Extending the same search cannot remove this obstruction; a new necessary
constraint or a stronger mathematical reduction is needed.

The full candidate universe has 61,884 groups and 641,524 integer systems.
The supplied partial certificate covers 3,651 groups and 95,443 integer roots,
leaving 58,233 groups outside this certificate. It uses 387 integer weight
vectors, 3,703 symbolic intervals and 38 binary nodes; its 95,481 final leaf
systems have minimum normalized strict gap $1/46$. There are 3,354 coarse
and 297 nested groups, 30 tree roots, and maximum tree depth two.
Every child of each included binary split is verified. The uncovered root
is explicitly required to be absent from the partial certificate.

The scalar audit checks 45 main lattice orientations, 16 intrinsic repairs,
50 early retention cases with minimum strict surplus 135, and 5,876 later
side/state prefix budgets. The two high-group row audits check 2,274,424
implications. Eleven negative controls include treating a partial certificate
as complete, omitting a witness branch, choosing a wrong branch, and altering
a witness coordinate. A separate program evaluates the emitted geometric
rows on 150 actual finite sets, checking 48,734 root-pair, intersection,
absorption, outer-pair, quantile and disjoint-allocation statements.

\paragraph{Execution record.}
The standard and \texttt{-O} progress outputs are byte-identical.
The complete numerical expansion of the supplied partial certificate is
recorded separately in eight batches; it concerns only the 95,443 supplied
integer roots, not the entire candidate universe. Final execution receipts,
source hashes and compressed outputs are stored in
\path{results/2026-10-18/upper88_progress/}.
The upper-89 foundation replay is byte-identical to its previous archive.
The archived upper bound 84 has not been independently replayed in full.
The general arguments above are written mathematical proofs; finite checks
supplement their review and are not proof-assistant formalization.

An earlier exploratory model placed the hull-width coefficients on the
wrong side of the pair inequality. That search was withdrawn before any
certificate or upper-bound claim was published. The corrected signs are
displayed above and checked directly on actual-set matrix assignments.
The unsuccessful searches are not counted as successful proof runs.

\paragraph{Reproduction.}
From the repository root run
{\footnotesize
\begin{verbatim}
python -I -S -B verification/upper88_progress.py --output .work/progress.json
python -I -S -B experiments/verify_upper88_geometric_rows.py --output .work/rows.json
\end{verbatim}
}
The first program reports \texttt{PASS\_PROGRESS\_ONLY}.
It never asserts $N\le88$. Windows users may run
\path{verify_upper88_progress_windows.bat}.
The optional \path{experiments/search_upper88_candidate.py} needs NumPy and
SciPy; all acceptance checks use only Python's standard library.
Compile this report twice with XeLaTeX.

\begin{thebibliography}{9}
\bibitem{previous} Project report, \emph{Compatible defects, retained prefixes,
and the audited bound 89}, \path{upper89_from90_20261018.tex}.
\bibitem{schedule} Project report, \emph{Variable-threshold persistent holes},
\path{upper90_from91_20261018.tex}, Sections 2--3.
\bibitem{foundation} Project report, \emph{From 115 to 103},
\path{upper103_from115_20261014.tex}.
\bibitem{archive} Archived structural notes,
\path{verification/upper84/progress85/inverse/proof85_inverse.md},
Sections 3--4. Method source only; endpoint conclusion not imported.
\bibitem{stanchescu} Y. Stanchescu, \emph{On addition of two distinct sets of
integers}, Acta Arithmetica \textbf{75} (1996), 191--194,
Theorem 4 and Corollary 2. \url{https://doi.org/10.4064/aa-75-2-191-194}.
\end{thebibliography}
\end{document}
